% ══════════════════════════════════════════════════════════════════════════════
%  CHAPTER 5: NUCLEAR SYNTHESIS IN STARS
% ══════════════════════════════════════════════════════════════════════════════
\chapter{Nuclear Synthesis in Stars}

% ═════════════════════════════════════════════════════════════════════════════ 

Recall that there are three primary energy sources in stars: gravitational contraction, radioactive decay and nuclear fusion. 

In this chapter, we will focus on the process of thermonuclear fusion, which is the viable long-term stellar energy source in most stars.

\section{Requisites for Thermonuclear Fusion}

\subsection{Nuclear Size and Binding Energy}
\begin{theory}
\begin{align}
    R &= R_0A^{1/3}, \qquad R_0\approx1.2\times10^{-15}\,\text{m}, \\
    BE(Z,N) &= \left[Zm_p+Nm_n-M(Z,N)\right]c^2.
\end{align}
Binding energy per nucleon increases for light nuclei and decreases for very heavy nuclei due to Coulomb repulsion.
\end{theory}

\begin{example}{}{ex:binding}
Calculate the binding energy per nucleon for helium-4 ($^4$He) and $^{12}$C.

\smallskip\textbf{Solution.}
\begin{align*} 
BE(^4\text{He}) &= \left[2m_p + 2m_n - M(^4\text{He})\right]c^2 \approx 28.3\,\text{MeV}, \\
BE(^4\text{He})/4 &\approx 7.1\,\text{MeV/nucleon}, \\
BE(^{12}\text{C}) &= \left[6m_p + 6 m_n - M(^{12}\text{C})\right]c^2 \approx 92.2\,\text{MeV}, \\
BE(^{12}\text{C})/12 &\approx 7.7\,\text{MeV/nucleon}.
\end{align*}   
\end{example}

\begin{theory}
The $Q$-value of a nuclear reaction is the net energy released or absorbed during the reaction, calculated as the difference in total mass-energy between the reactants and products. For a reaction $A + B \rightarrow C + D$, the $Q$-value is given by:
\begin{equation}
    Q = \left[m_A + m_B - (m_C + m_D)\right]c^2.    
\end{equation}
A positive $Q$-value indicates an exothermic reaction, while a negative $Q$-value indicates an endothermic reaction.
\end{theory}

\subsection{Why Fusion Powers Stars}
\begin{intuition}
Stars gain energy when fusing light nuclei toward higher binding energy per nucleon.
This trend reverses near iron, where further fusion is endothermic.
\end{intuition}

\section{Hydrogen and Helium Burning}
\subsection{Proton-Proton (p-p) Chain}

\begin{foundation}
Representative net reaction:
\begin{equation}
    4\,{}^1\text{H}\rightarrow{}^4\text{He}+2e^++2\nu_e+\text{energy}.
\end{equation}
\end{foundation}

\begin{example}{}{ex:ppchain}
Calculate the energy released in the proton-proton chain reaction.

\smallskip\textbf{Solution.}
The net reaction for the proton-proton chain is:    
\begin{equation}
    4\,{}^1\text{H}\rightarrow{}^4\text{He}+2e^++2\nu_e+\text{energy}.  
\end{equation}
The energy released can be calculated using the mass difference between the reactants and products:
\begin{align*}
\Delta m &= 4m_p - m_{^4\text{He}} - 2m_e \\
&\approx 4 \times 1.007276\,\text{u} - 4.002603\,\text{u} - 2 \times 0.0005486\,\text{u} \\
&\approx 0.0287\,\text{u}.  
\end{align*}
Converting this mass difference to energy using $E = \Delta m c^2$ and
knowing that $1\,\text{u} \approx 931.5\,\text{MeV}/c^2$, we get:
\begin{align*}
E &= 0.0287\,\text{u} \times 931.5\,\text{MeV}/\text{u} \\
&\approx 26.7\,\text{MeV}.
\end{align*}
Thus, the energy released in the proton-proton chain reaction is approximately $26.7\,\text{MeV}$ and carried away by neutrinos.
\end{example}

\subsection{Helium Burning (Triple-Alpha Process)}
\begin{foundation}
Representative net reaction:
\begin{equation}
    3\,{}^4\text{He}\rightarrow{}^{12}\text{C}+\gamma+\text{energy}.
\end{equation}
\end{foundation}
\begin{example}{}{ex:triplealpha}
Calculate the energy released in the triple-alpha process.

\smallskip\textbf{Solution.}
The net reaction for the triple-alpha process is:    
\begin{equation}
    3\,{}^4\text{He}\rightarrow{}^{12}\text{C}+\gamma+\text{energy}.  
\end{equation}
The energy released can be calculated using the mass difference between the reactants and products:
\begin{align*}
\Delta m &= 3m_{^4\text{He}} - m_{^{12}\text{C}} \\
&\approx 3 \times 4.002603\,\text{u} - 12.000000\,\text{u} \\
&\approx 0.007809\,\text{u}.  
\end{align*}
Converting this mass difference to energy using $E = \Delta m c^2$ and
knowing that $1\,\text{u} \approx 931.5\,\text{MeV}/c^2$, we get:
\begin{align*}
E &= 0.007809\,\text{u} \times 931.5\,\text{MeV}/\text{u} \\
&\approx 7.27\,\text{MeV}.
\end{align*}
Thus, the energy released in the triple-alpha process is approximately $7.27\,\text{MeV}$.
\end{example}

\section{CNO Cycle}
The CNO cycle is a catalytic cycle that uses carbon, nitrogen, and oxygen as catalysts to convert hydrogen into helium. 
It dominates in stars more massive than the Sun due to its stronger temperature dependence compared to the p-p chain.
\begin{foundation}
Representative net reaction:
\begin{align}
^12\text{C} + p &\rightarrow {}^{13}\text{N} + \gamma, \\
{}^{13}\text{N} &\rightarrow {}^{13}\text{C} + e^+ + \nu_e, \\
{}^{13}\text{C} + p &\rightarrow {}^{14}\text{N} + \gamma, \\
{}^{14}\text{N} + p &\rightarrow {}^{15}\text{O} + \gamma, \\
{}^{15}\text{O} &\rightarrow {}^{15}\text{N} + e^+ + \nu_e, \\
{}^{15}\text{N} + p &\rightarrow {}^{12}\text{C} + {}^4\text{He}.
\end{align}
\end{foundation}

\section{Thermonuclear Reaction Rates}

\subsection{Coulomb Barrier and Tunneling}

\begin{foundation}
The \textbf{Gamov energy} is the minimum kinetic energy that two nuclei must have in order to overcome the Coulomb barrier and undergo fusion. It is given by the formula:
\begin{equation}    
E_G = (\pi \alpha Z_1 Z_2)^2 2 m_{r}c^2
\end{equation} 
where $m_r$ is the reduced mass of the two nuclei, $\alpha$ is the fine-structure constant, and $Z_1$ and $Z_2$ are the atomic numbers of the two nuclei.
\end{foundation}
\begin{example}{The fusion of two protons}{ex:gamov1}
Calculate the Gamov energy $E_G$ for the fusion of two protons.

\smallskip\textbf{Solution.}
For two protons, we have $Z_1 = Z_2 = 1$ and the reduced mass $m_r$ is approximately equal to the mass of a proton, $m_p$. The fine-structure constant $\alpha$ is approximately $1/137$. Plugging these values into the formula for the Gamov energy, we get:
\begin{align*}
E_G &= (\pi \alpha Z_1 Z_2)^2 2 m_{r}c^2 \\
&= (\pi \times \frac{1}{137} \times 1 \times 1)^2 \times 2 m_p c^2 \\
&= \left(\frac{\pi}{137}\right)^2 \times 2 m_p c^2 \\
&\approx 493\,\text{keV}.
\end{align*}
Thus, the Gamov energy for the fusion of two protons is approximately $493\,\text{keV}$.
\end{example}

\begin{example}{The fusion of carbon and oxygen}{ex:gamov2}
Calculate the Gamov energy $E_G$ for the fusion of carbon-12 ($^{12}\text{C}$) and oxygen-16 ($^{16}\text{O}$). 

\smallskip\textbf{Solution.}
For carbon-12, we have $Z_1 = 6$ and for oxygen-16, we have $Z_2 = 8$. The reduced mass $m_r$ can be calculated using the formula:
\begin{equation}
m_r = \frac{m_1 m_2}{m_1 + m_2},
\end{equation}
where $m_1$ and $m_2$ are the masses of the carbon and oxygen nuclei, respectively. The mass of carbon-12 is approximately $12 m_p$ and the mass of oxygen-16 is approximately $16 m_p$. Thus, we have:
\begin{align*}
m_r &= \frac{(12 m_p)(16 m_p)}{12 m_p + 16 m_p} \\
&= \frac{192 m_p^2}{28 m_p} \\        
&= \frac{192}{28} m_p \\
&\approx 6.86 m_p.
\end{align*}
Now we can calculate the Gamov energy using the formula:
\begin{align*}
E_G &= (\pi \alpha Z_1 Z_2)^2 2 m_{r}c^2 \\
&= (\pi \times \frac{1}{137} \times 6 \times 8)^2 \times 2 \times 6.86 m_p c^2 \\
&= \left(\frac{48\pi}{137}\right)^2 \times 2 \times 6.86 m_p c^2 \\
&\approx 8.7\,\text{MeV}. 
\end{align*}
Thus, the Gamov energy for the fusion of carbon-12 and oxygen-16 is approximately $8.7\,\text{MeV}$.
\end{example}


\begin{foundation}
The electrostatic interaction (Coulomb barrier) between two positively charged nuclei is given by:
\begin{equation}
V_C(r) = \frac{Z_1 Z_2 e^2}{4\pi \epsilon_0 r}.
\end{equation}
Classically, nuclei with $E<V_C$ cannot fuse as they do not have enough kinetic energy to overcome the Coulomb barrier.

Quantum tunneling gives non-zero penetration probability $P$ for $E<V_C$. The probability of tunneling through the Coulomb barrier can be approximated by the formula:
\begin{equation}
    P \propto \exp\left[-\sqrt{\frac{E_G}{E}}\right],
\end{equation}
where $E_G$ is the Gamow energy.
\end{foundation}

\section{Thermonuclear Reaction Rates}

\begin{theory}
The cross-section $\sigma v$ for a nuclear reaction can be expressed as:
\begin{equation}
    \langle\sigma v\rangle = \left(\frac{8}{\pi\mu}\right)^{1/2}
    \frac{1}{(k_BT)^{3/2}}\int_0^\infty S(E)
    \exp\left[-\frac{E}{k_BT}-\sqrt{\frac{E_G}{E}}\right]dE.
\end{equation}
The integrand peaks at the Gamow energy window, balancing high-energy tail abundance and tunneling probability.
\end{theory}

\begin{example}{}{ex:reactionrate}
Find the peak of the integrand in the expression for $\langle\sigma v\rangle$ to determine the most effective energy for nuclear reactions at a given temperature.

\smallskip\textbf{Solution.}
To find the peak of the integrand, we need to maximize the function:
\begin{equation}
f(E) = \exp\left[-\frac{E}{k_BT}-\sqrt{\frac{E_G}{E}}\right].
\end{equation}
Taking the natural logarithm of $f(E)$, we get:
\begin{equation}
\ln f(E) = -\frac{E}{k_BT} - \sqrt{\frac{E_G}{E}}.
\end{equation}  
To find the maximum, we take the derivative with respect to $E$ and set it to zero:
\begin{equation}
\frac{d}{dE} \ln f(E) = -\frac{1}{k_BT} + \frac{1}{2} \sqrt{\frac{E_G}{E^3}} = 0.
\end{equation}
Rearranging this equation gives:
\begin{equation}
\frac{1}{k_BT} = \frac{1}{2} \sqrt{\frac{E_G}{E^3}}.
\end{equation}
Squaring both sides, we get:
\begin{equation}
\frac{1}{(k_BT)^2} = \frac{1}{4} \frac{E_G}{E^3}.
\end{equation}
Rearranging this equation gives:
\begin{equation}
E^3 = \frac{1}{4} E_G (k_BT)^2.
\end{equation}  
Taking the cube root of both sides, we find the most effective energy for nuclear reactions at a given temperature:
\begin{equation}
E_0 = \left(\frac{1}{4} E_G (k_BT)^
2}\right)^{1/3} = \left(\frac{E_G}{4}\right)^{1/3} (k_BT)^{2/3}.
\end{equation}
Thus, the most effective energy for nuclear reactions at a given temperature is approximately $E_0 = \left(\frac{E_G}{4}\right)^{1/3} (k_BT)^{2/3}$.
\end{example}


\begin{theory}
\begin{equation}
    R_{AB} \propto T^{a}
\end{equation}
where $a = \sqrt[3]{\frac{E_G}{4kT}}$ is a positive exponent that depends on the specific reaction and the temperature range.
\end{theory}

\begin{example}{}{ex:ppchain}
Calculate the temperature dependence of the proton-proton (p-p) chain reaction rate.

\smallskip\textbf{Solution.}
The proton-proton chain reaction has a Gamow energy of approximately $E_G \approx 493\,\text{keV}$. The temperature dependence can be approximated by:
\begin{equation}
    R_{pp} \propto T^{\sqrt[3]{\frac{E_G}{4kT}}}.           
\end{equation}
At typical stellar core temperatures of around $T \approx 1.5 \times 10^7\,\text{K}$, the exponent $a$ can be calculated as:
\begin{equation}
    a = \sqrt[3]{\frac{493\,\text{keV}}{4k(1.5 \times 10^7\,\text{K})}} \approx 4.5.    
\end{equation}
Thus, the reaction rate for the proton-proton chain scales approximately as $R_{pp} \propto T^{4.5}$ at typical stellar core temperatures.
\end{example}